\subsection{Bidual}

DEFINITION 1. --- Let E be a locally convex space and \(E_{b}^{\prime}\) its strong dual. The dual of the locally convex space \(E_{b}^{\prime}\) is called the bidual of E and is denoted by \(E^{\prime\prime}\) .

For every \(x \in E\) , let \(\tilde{x}\) be the linear form \(x' \mapsto \langle x, x' \rangle\) on \(E'\) ; it is continuous for the weak topology \(\sigma(E', E)\) , hence a fortiori, for the strong topology on \(E'\) ; therefore \(\tilde{x} \in E''\) for all \(x \in E\) . The map \(c_{E}: x \mapsto \tilde{x}\) from E into \(E''\) is a linear mapping, said to be canonical.

PROPOSITION 1. --- The kernel of \(c_{\mathrm{E}}: \mathrm{E} \to \mathrm{E}''\) is the closure of 0 in E. If E is Hausdorff, \(c_{\mathrm{E}}\) is injective.

By construction, the kernel of \(c_{\mathrm{E}}\) is the intersection of the kernels of the continuous linear forms on E, i.e.~the closure of \(\{0\}\) in E (II, p.~24, cor. 1).

When E is Hausdorff, we identify E with a subspace of \(E''\) , by means of the mapping \(c_{E}\) .

The strong topology on \(E''\) is the S-topology, where S is the family of all strongly bounded subsets of \(E'\) . Since every equicontinuous subset of \(E'\) is strongly bounded (III, p.~22, prop. 9), the initial topology on E is coarser than the topology obtained by taking the inverse image under \(c_{E}\) of the strong topology on \(E''\) ; it can be strictly coarser (IV, p.~52, exerc. 1). However :

PROPOSITION 2. --- Suppose that the space E is bornological or barrelled. The initial topological on E is the inverse image under \(c_{E}\) of the strong topology on \(E''\) .

For, every subset of E' which is strongly bounded is equicontinuous (III, p.~22, prop. 10 and III, p.~24).

PROPOSITION 3. --- Let E be a locally convex Hausdorff space. In order that the strong dual \(E_{b}^{\prime}\) of E be barrelled, it is necessary and sufficient that every subset of \(E^{\prime\prime}\) which is bounded for \(\sigma(E^{\prime\prime}, E^{\prime})\) , is contained in the closure, for \(\sigma(E^{\prime\prime}, E^{\prime})\) , of a bounded subset of E.

The equicontinuous subsets of \(E''\) are the subsets contained in the bipolar (for the duality between \(E''\) and \(E'\) ) of a bounded subset of the subspace E of \(E''\) . It is now enough to apply the theorem of bipolars (II, p.~45, cor. 3) and the definition of a barrelled space (III, p.~24).

Remark. --- Let E be a locally convex Hausdorff space, \(E'\) its dual and \(E''\) its bidual. We have \(E \subset E'' \subset E'^{*}\) , where \(E'^{*}\) is the algebraic dual of \(E'\) . If B is a bounded subset of E, its closure \(\overline{B}\) in \(E'^{*}\) endowed with \(\sigma(E'^{*}, E')\) is contained in \(E''\) : for, the polar \(U = B^{\circ}\) of B in \(E'\) is a neighbourhood of 0 in \(E_b'\) , and we have \(\overline{B} \subset U^{\circ} \subset E''\) .

